\section{Indledning}

\emph{Digitale træningsapps} kan gøre træningsprogrammer, øvelser, historik og feedback lettere tilgængelige. Det løser dog kun en del af problemet. Vedvarende fysisk aktivitet hænger også sammen med autonom motivation, oplevet kompetence og støtte, der giver brugeren en meningsfuld oplevelse af kontrol \cite{teixeira2012_sdt_exercise}. Digital brug har først værdi, når den faktisk støtter den ønskede adfærd og ikke blot skaber mere skærmtid \cite{yardley2016_effective_engagement_dbci}.

I praksis møder træningsplanen en hverdag med begrænset tid, ændret dagsform, fyldt træningscenter, skiftede øvelser og ufuldstændig \emph{logging}. Den slags besvær betegnes som \emph{friktion}. Hvis appen kun giver mening ved perfekt efterlevelse, bliver støtten svag netop i de situationer, hvor brugeren har mest brug for overblik.

Den afgørende forskel ligger mellem en plan, der ser rigtig ud på forhånd, og et forløb, der stadig giver mening bagefter. En kalender kan fortælle, hvad brugeren skulle have gjort. En brugbar træningsapp skal også kunne håndtere, hvad brugeren faktisk gjorde, og omsætte det til en næste handling uden at gøre ændringen til et nederlag.

Apps og aktivitetsmålere kan understøtte fysisk aktivitet, men studierne peger samtidig på, at effekten varierer med design, brugskontekst og vedvarende anvendelse \cite{laranjo2021_apps_trackers_meta,romeo2019_smartphone_apps_pa_meta,brickwood2019_wearable_trackers_meta}. Frafald er derfor et centralt vilkår i digitale sundhedsinterventioner. Når brugeren stopper, før interventionen kan få virkning, mister løsningen sin praktiske værdi \cite{eysenbach2005_law_of_attrition}.

I styrketræning bliver udfordringen mere konkret, fordi et program rummer mere end en kalender. American College of Sports Medicine (ACSM) beskriver progression gennem belastning, volumen, øvelsesvalg, frekvens og løbende justering \cite{acsm2009_progression_models}. Nyere resistance-training-anbefalinger fastholder progression som et centralt princip, mens forskning i autoregulering gør justering efter dagsform og performance fagligt relevant \cite{currier2026_acsm_resistance_training,greig2020_autoregulation_resistance_training,larsen2021_autoregulation_systematic_review}.

TræningsMester er udviklet som en \emph{native iOS-app}, hvor denne spænding bliver konkret. Plan, tracker, completion, feedback og næste handling mødes i samme brugerforløb. Brugeren kan springe et pas over, ændre belastning, bytte øvelser, træne uden at logge alle sæt eller vende tilbage efter en afbrudt session. Appens centrale udfordring er at fastholde retning og progression, når træningen ikke følger en perfekt lineær plan.

\tmdelkonklusion{Mellempointe}{Problemet er bruddet mellem plan og hverdag. En brugbar træningsapp skal kunne vise, at brugeren stadig er i gang, og hvad næste skridt er, også når træningen ikke blev logget perfekt.}
